% Part: many-valued-logic
% Chapter: three-valued-logics
% Section: goedel

\documentclass[../../../include/open-logic-section]{subfiles}

\begin{document}

\olfileid[id]{mvl}{thr}{god}

\olsection{Logika G\"odel}

Kurt G\"odel memperkenalkan suatu barisan logika bernilai~$n$ yang masing-masing
memuat semua !!{formula}s yang valid dalam logika intuisionistik dan termuat
dalam logika klasik. Berikut adalah logika pertama yang menarik dalam barisan
tersebut:

\begin{defn}\ollabel{defn:goedel}
\emph{Logika G\"odel bernilai~$3$}~$\LogGod$ didefinisikan dengan menggunakan
matriks berikut:
\begin{enumerate}
  \item Bahasa proposisional baku~$\Lang L_0$ dengan $\lfalse$, $\lnot$,
  $\land$, $\lor$, $\lif$.
  \item Himpunan nilai kebenaran~$V=\{\True,\Undef,\False\}$.
  \item $\True$ merupakan satu-satunya nilai yang ditetapkan, yakni
  $V^+=\{\True\}$.
  \item Bagi $\lfalse$, kita mempunyai $\tf{\lfalse}=\False$. Fungsi
  kebenaran bagi konektif lainnya diberikan oleh tabel-tabel berikut:
  \begin{center}
    \begin{tabular}{c|c}
      $\tf{\lnot}[\LogGod]$ & \\
      \hline
      $\True$ & $\False$ \\
      $\Undef$ & $\False$ \\
      $\False$ & $\True$
    \end{tabular}
    \quad
    \begin{tabular}{c|ccc}
      $\tf{\land}[\LogGod]$ & $\True$ & $\Undef$ & $\False$ \\
      \hline
      $\True$ & $\True$ & $\Undef$ & $\False$ \\
      $\Undef$ & $\Undef$ & $\Undef$ & $\False$\\
      $\False$ & $\False$ & $\False$ & $\False$
    \end{tabular}
    \\[2ex]
    \begin{tabular}{c|ccc}
      $\tf{\lor}[\LogGod]$ & $\True$ & $\Undef$ & $\False$ \\
      \hline
      $\True$ & $\True$ & $\True$ & $\True$ \\
      $\Undef$ & $\True$ & $\Undef$ & $\Undef$ \\
      $\False$ & $\True$ & $\Undef$ & $\False$
    \end{tabular}
    \quad
    \begin{tabular}{c|ccc}
      $\tf{\lif}[\LogGod]$ & $\True$ & $\Undef$ & $\False$ \\
      \hline
      $\True$ & $\True$ & $\Undef$ & $\False$ \\
      $\Undef$ & $\True$ & $\True$ & $\False$  \\
      $\False$ & $\True$ & $\True$ & $\True$
    \end{tabular}
  \end{center}
\end{enumerate}
\end{defn}

Perhatikan bahwa tabel kebenaran bagi $\land$ dan~$\lor$ sama seperti dalam
logika \L ukasiewicz dan logika Kleene kuat, sedangkan tabel kebenaran bagi
$\lnot$ dan~$\lif$ berbeda pada masing-masing logika. Dalam logika G\"odel,
$\tf{\lnot}(\Undef)=\False$. Berbeda dari logika \L ukasiewicz dan logika
Kleene, $\tf{\lif}(\Undef,\False)=\False$; berbeda dari logika Kleene
(tetapi sama seperti dalam logika \L ukasiewicz),
$\tf{\lif}(\Undef,\Undef)=\True$.

Sebagaimana disiratkan oleh hubungan dengan logika intuisionistik yang disebut
di atas, $\LogGod[3]$ dekat dengan logika intuisionistik. Semua kebenaran
intuisionistik merupakan tautologi dalam~$\LogGod[3]$, dan banyak tautologi
klasik yang tidak valid secara intuisionistik juga bukan tautologi
dalam~$\LogGod[3]$. Sebagai contoh, formula-formula berikut bukan tautologi:
\begin{align*}
  & p \lor \lnot p && (p \lif q) \lif (\lnot p \lor q) \\
  & \lnot\lnot p \lif p && \lnot(\lnot p \land \lnot q) \lif (p \lor q) \\
  & ((p \lif q) \lif p) \lif p && \lnot(p \lif q) \lif (p \land \lnot q)
\end{align*}
Akan tetapi, tidak setiap tautologi dari~$\LogGod[3]$ valid secara
intuisionistik; contohnya ialah $\lnot\lnot p\lor\lnot p$ dan
$(p\lif q)\lor(q\lif p)$.

\begin{prob}
  Berikan tabel kebenaran untuk menunjukkan bahwa formula-formula berikut
  merupakan tautologi dari~$\LogGod[3]$:
  \begin{align*}
    & \lnot\lnot p \lor \lnot p\\
    & (p \lif q) \lor (q \lif p) \\
    & \lnot(p \land q) \lif (\lnot p \lor \lnot q) \\
    & (p \lif q) \lor (q \lif r) \lor (r \lif s)
  \end{align*}
\end{prob}

\begin{prob}
  Berikan tabel kebenaran untuk menunjukkan bahwa formula-formula berikut
  bukan tautologi dari~$\LogGod[3]$:
  \begin{align*}
    & (p \lif q) \lif (\lnot p \lor q) \\
    & \lnot(\lnot p \land \lnot q) \lif (p \lor q) \\
    & ((p \lif q) \lif p) \lif p \\
    & \lnot(p \lif q) \lif (p \land \lnot q)
  \end{align*}
\end{prob}

\begin{prob}
  Manakah di antara relasi-relasi berikut yang berlaku dalam logika G\"odel?
  Berikan tabel kebenaran bagi masing-masing relasi.
  \begin{enumerate}
    \item $p,p\lif q\Entails q$
    \item $p\lor q,\lnot p\Entails q$
    \item $p\land q\Entails p$
    \item $p\Entails p\land p$
    \item $p\Entails p\lor q$
  \end{enumerate}
\end{prob}

\end{document}
